%%% FIELD proof-certified-gaussian-value
Put \(I=I(\theta_0)\), \(q_{\mathrm{cov}}=1-\alpha\), \(P_u=N(u,I^{-1})\), and
\[
\mathcal L(C)=\sup_{t\in\mathcal D}\operatorname{diam}C(Z,t).
\]
The covariance spectrum assumption and the positive definiteness in the definition of \(I(\theta_0)\) make every inverse and whitening operation below legitimate.  By \(\mathrm{lem:local\mbox{-}risk\mbox{-}map}\), \(\Psi_{\theta_0}\) is continuous and affine on each member of a finite polyhedral complex on the compact set \(\mathcal U\times\mathcal D\).  Let \(K_u\) be the largest operator norm, over those affine pieces, of the coefficient multiplying \(u\).  Partitioning any line segment from \(u\) to \(u'\) at its finitely many cell crossings and summing its affine increments gives
\[
\sup_{t\in\mathcal D}
\left|\Psi_{\theta_0}(u,t)-\Psi_{\theta_0}(u',t)\right|
\le K_u\lVert u-u'\rVert_2.
\]
Thus
\[
D_\star=\sup_{u,u'\in\mathcal U}d_\Psi(u,u')
\le 2K_uM<\infty.
\]

We first make both approximation errors explicit.  Whitening two Gaussian laws with the same covariance and integrating in the single likelihood-ratio direction gives
\[
\operatorname{TV}(P_u,P_{u_i})
=2\Phi\!\left(\frac{\lVert I^{1/2}(u-u_i)\rVert_2}{2}\right)-1.
\]
The derivative of the right-hand side, as a function of the displayed norm, is at most \(1/\sqrt{2\pi}\).  Hence, if \(\lVert u-u_i\rVert_2\le\delta\),
\[
\operatorname{TV}(P_u,P_{u_i})
\le \frac{\lVert I^{1/2}\rVert_{\mathrm{op}}\delta}{\sqrt{2\pi}}
=:\epsilon_\delta.
\]

For observation quantization, write \(P_0=N(0,I^{-1})\).  On each bounded cell \(o\), reconstruct from the parameter-independent conditional distribution \(P_0(\,\cdot\mid o)\), and use any fixed probability distribution on the tail cell.  The likelihood ratio is
\[
\ell_u(z)=\frac{dP_u}{dP_0}(z)
=\exp\!\left\{u^{\prime}Iz-\frac12u^{\prime}Iu\right\}.
\]
If \(o\) has diameter at most \(h\), then
\[
\sup_{z,z'\in o}
\left|\log\ell_u(z)-\log\ell_u(z')\right|
\le \lVert I\rVert_{\mathrm{op}}Mh.
\]
Conditionally on \(o\), the density ratio between \(P_u(\,\cdot\mid o)\) and the reconstruction law is \(\ell_u(z)/E_0[\ell_u(Z)\mid Z\in o]\).  The preceding oscillation bound therefore makes their total variation distance no larger than
\(\exp\{\lVert I\rVert_{\mathrm{op}}Mh\}-1\).  If \(R>M\), a coordinatewise Gaussian tail bound gives
\[
\tau_R:=\sup_{\lVert u\rVert_2\le M}
P_u\{Z\notin[-R,R]^q\}
\le 2q\exp\!\left\{
-\frac{(R-M)^2}{2\lambda_{\max}(I^{-1})}
\right\}.
\]
Consequently, if \(\widetilde P_u\) is the law obtained by recording the cell and applying the reconstruction kernel,
\[
\sup_{u\in\mathcal U}\operatorname{TV}(P_u,\widetilde P_u)
\le
\epsilon_h:=
\min\!\left\{1,\,
\exp\{\lVert I\rVert_{\mathrm{op}}Mh\}-1+\tau_R
\right\}.
\]
This bound tends to zero when \(R\to\infty\) and \(h\to0\); in particular it does so under \(R\to\infty\) and \(hR\to0\).

We next prove the lower certificate.  Let \(C\in\mathfrak C_{1-\alpha}\) be any externally randomized measurable interval-section band.  Apply \(\mathrm{prop:curve\mbox{-}decision\mbox{-}reduction}\) to \(F=\mathcal U_N\), and let \(A\) be the random subset of net points whose complete curves are covered.  Pointwise,
\[
c_A\le\mathcal L(C),
\qquad
P_{u_i}\{i\in A\}\ge q_{\mathrm{cov}}.
\]
Quantize, reconstruct, and then run this acceptance rule, including its original external randomization.  Conditional on the observed cell \(o\), let \(x_{oA}\) be the probability of selecting \(A\).  Then \(x_{oA}\ge0\) and \(\sum_Ax_{oA}=1\).  Total variation changes the probability of any event by at most \(\epsilon_h\), while \(0\le c_A\le D_\star\); hence, for every \(i\),
\[
\sum_op_{io}\sum_{A\ni i}x_{oA}
\ge q_{\mathrm{cov}}-\epsilon_h
\]
and
\[
\sum_op_{io}\sum_Ac_Ax_{oA}
\le E_{u_i}[\mathcal L(C)]+D_\star\epsilon_h.
\]
Thus \(x\) is feasible for the finite program at coverage \(q_{\mathrm{cov}}-\epsilon_h\), with objective at most
\(\sup_uE_u[\mathcal L(C)]+D_\star\epsilon_h\).  Taking the infimum over \(C\) proves
\[
V_{N,h}(q_{\mathrm{cov}}-\epsilon_h)-D_\star\epsilon_h
\le L_\alpha(\theta_0,v_0;M,L_0).
\]

For the upper certificate, take a finite-program rule feasible at coverage \(q_{\mathrm{cov}}+\epsilon_\delta\), with objective no larger than
\(V_{N,h}(q_{\mathrm{cov}}+\epsilon_\delta)+\eta\), where \(\eta>0\).  If the rule selects a nonempty set \(A\), output
\[
C_A^\delta(t)=
\left[
\min_{i\in A}\Psi_{\theta_0}(u_i,t)-K_u\delta,\,
\max_{i\in A}\Psi_{\theta_0}(u_i,t)+K_u\delta
\right];
\]
if \(A=\varnothing\), output a fixed singleton.  Given \(u\in\mathcal U\), choose a nearest \(u_i\), breaking ties by a fixed ordering, so that \(\lVert u-u_i\rVert_2\le\delta\).  On \(\{i\in A\}\), the Lipschitz inequality proves simultaneous coverage of the entire curve indexed by \(u\).  Therefore
\[
\begin{aligned}
P_u\{\Psi_{\theta_0}(u,t)\in C_A^\delta(t)
\text{ for all }t\in\mathcal D\}
&\ge P_{u_i}\{i\in A\}
-\operatorname{TV}(P_u,P_{u_i})\\
&\ge q_{\mathrm{cov}}.
\end{aligned}
\]
The constructed band is measurable because the partition and parameter net are finite, and it belongs to \(\mathfrak C_{1-\alpha}\).  Its loss is at most \(c_A+2K_u\delta\).  A second total-variation comparison yields
\[
\sup_{u\in\mathcal U}E_u[\mathcal L(C_A^\delta)]
\le
V_{N,h}(q_{\mathrm{cov}}+\epsilon_\delta)
+\eta+D_\star\epsilon_\delta+2K_u\delta.
\]
Letting \(\eta\downarrow0\) proves the stated upper certificate.

We now verify finite-program duality.  Attach nonnegative multipliers \(\ell_i\) to the coverage constraints, nonnegative multipliers \(\pi_i\) to the risk constraints, and unrestricted multipliers \(\eta_o\) to the row-normalization constraints.  The Lagrangian is
\[
\begin{aligned}
&\zeta\left(1-\sum_i\pi_i\right)
+q\sum_i\ell_i+\sum_o\eta_o\\
&\quad+
\sum_{o,A}x_{oA}
\left(
c_A\sum_i\pi_ip_{io}
-\sum_{i\in A}\ell_ip_{io}
-\eta_o
\right),
\end{aligned}
\]
where \(q\) denotes the coverage argument of the finite program.  Its infimum over unrestricted \(\zeta\) and \(x_{oA}\ge0\) is finite exactly if
\[
\sum_i\pi_i=1,
\qquad
\eta_o\le
c_A\sum_i\pi_ip_{io}
-\sum_{i\in A}\ell_ip_{io}
\quad\text{for every }(o,A),
\]
which are precisely the displayed dual constraints.  This gives weak duality.

For the reverse inequality, fix \(\gamma<V_{N,h}(q)\).  Let \(\mathcal X\) be the product of the finitely many subset-probability simplices and define
\[
a_i(x)=\sum_op_{io}\sum_{A\ni i}x_{oA},
\qquad
b_i(x)=\sum_op_{io}\sum_Ac_Ax_{oA}.
\]
The compact convex set
\[
\mathcal K=
\left\{
\bigl((q-a_i(x))_i,(b_i(x)-\gamma)_i\bigr):
x\in\mathcal X
\right\}
\]
is disjoint from the nonpositive orthant.  Let a closest pair realize the strictly positive distance between these two closed convex sets.  Expanding the squared-distance inequality after moving the point of \(\mathcal K\) toward any other point of \(\mathcal K\) gives a separating normal with nonnegative coordinates.  Hence there are \(\ell_i,\pi_i\ge0\), not all zero, such that
\[
\sum_i\ell_i\{q-a_i(x)\}
+\sum_i\pi_i\{b_i(x)-\gamma\}>0
\quad\text{for every }x\in\mathcal X.
\]
The action that always selects the full net has \(a_i=1\ge q\) for every \(i\); substituting it shows \(\sum_i\pi_i>0\).  Divide all multipliers by this positive sum, so \(\sum_i\pi_i=1\), and put
\[
\eta_o=
\min_A\left\{
c_A\sum_i\pi_ip_{io}
-\sum_{i\in A}\ell_ip_{io}
\right\}.
\]
Minimizing separately over every simplex row in the strict separating inequality shows
\[
q\sum_i\ell_i+\sum_o\eta_o>\gamma.
\]
Thus the displayed \((\pi,\ell,\eta)\) is dual feasible with value exceeding every \(\gamma<V_{N,h}(q)\).  Weak duality now proves equality of the finite primal and dual values.  Primal and dual feasible points therefore provide checkable certificates.

It remains to prove convergence of the two certificates.  The hypotheses
\(\epsilon_h\le q_{\mathrm{cov}}\) and \(\epsilon_\delta\le\alpha\) ensure that both program coverage levels lie in \([0,1]\).  Mix an optimal rule at coverage \(q_{\mathrm{cov}}-\epsilon_h\) with the rule that always selects the full net.  The mixing probability
\[
\beta=
\frac{\epsilon_\delta+\epsilon_h}{\alpha+\epsilon_h}
\]
belongs to \([0,1]\), and the mixture has coverage at least
\[
(1-\beta)(q_{\mathrm{cov}}-\epsilon_h)+\beta
=q_{\mathrm{cov}}+\epsilon_\delta.
\]
Because every subset cost is at most \(D_\star\),
\[
V_{N,h}(q_{\mathrm{cov}}+\epsilon_\delta)
-V_{N,h}(q_{\mathrm{cov}}-\epsilon_h)
\le
D_\star\frac{\epsilon_\delta+\epsilon_h}{\alpha+\epsilon_h}.
\]
Combining this with the lower and upper certificates bounds their gap by
\[
D_\star\left\{
\frac{\epsilon_\delta+\epsilon_h}{\alpha+\epsilon_h}
+\epsilon_\delta+\epsilon_h
\right\}
+2K_u\delta,
\]
which tends to zero when \(\delta\to0\), \(R\to\infty\), and \(hR\to0\).

All inputs to these finite programs have certified finite computations.  By
\(\mathrm{prop:curve\mbox{-}decision\mbox{-}reduction}\), each \(c_A\) is obtained from finitely many second-order-cone programs; primal and conic-dual feasible points give outward numerical enclosures.  On \([-R,R]^q\), the Gaussian density and its derivatives are explicit and uniformly bounded for \(u_i\in\mathcal U\).  Interval quadrature on refined cells therefore encloses every \(p_{io}\), and \(\tau_R\) encloses the tail contribution.  Replacing costs and probabilities by outward endpoints and padding each coverage row by its rowwise probability error gives finite robust linear programs.  A rowwise probability error \(e_p\) changes a coverage row by at most \(e_p\) and a risk row by at most \(D_\star e_p\); a uniform cost error \(e_c\) changes a risk row by at most \(e_c\).  Hence the robust lower and upper numerical values converge to the exact finite-program value as \(e_p,e_c\to0\).

Finally, choose meshes and numerical enclosures satisfying the stated limits, and use each upper-program optimizer in the padded-envelope construction.  The vanishing analytic and numerical certificate gap proves that its worst expected maximum width converges to
\(L_\alpha(\theta_0,v_0;M,L_0)\).  Therefore these rules are asymptotically minimax among all externally randomized measurable interval-section bands.  If \(D_\star=0\), all curves coincide and the singleton band containing the common curve has coverage one and width zero, which proves the same conclusion directly.

%%% FIELD proof-panel-experiment-transfer
Index pre-treatment group \(\ell\) by dimension \(J+1\), count \(n_\ell\), mean
\[
m_\ell(\theta)=(a_{0\ell}(\theta),A_{\ell\cdot}(\theta))^{\prime},
\]
and covariance \(\Sigma(\theta)\).  Add the post-donor group of dimension \(J\), count \(T_1\), mean \(d(\theta)\), and covariance \(\Sigma_{DD}(\theta)\).  By the sampling-fraction assumption, every group count divided by \(n\) has a strictly positive finite limit.  By the covariance-spectrum assumption, all covariance inverses below exist uniformly on a neighborhood of \(\theta_0\).  The local-chart assumption gives twice continuously differentiable mean and covariance coordinates, and the Gaussian-submodel assumption gives independence across all displayed observations.

For a group \(g\), the directional score of one observation \(Y\sim N(\mu_g,\Sigma_g)\) in chart direction \(u\) is
\[
s_g(Y)[u]
=D\mu_g[u]^{\prime}\Sigma_g^{-1}(Y-\mu_g)
+\frac12\left\{
(Y-\mu_g)^{\prime}\Sigma_g^{-1}D\Sigma_g[u]
\Sigma_g^{-1}(Y-\mu_g)
-\operatorname{tr}(\Sigma_g^{-1}D\Sigma_g[u])
\right\}.
\]
Let \(\Delta_n\in\mathbb R^q\) be \(n^{-1/2}\) times the sum of the coordinate-score vectors over all groups.  Centered Gaussian odd moments make the linear mean-score part orthogonal to the centered quadratic covariance-score part.  The identities
\[
E[(a^{\prime}X)(b^{\prime}X)]=a^{\prime}\Sigma_gb
\]
and
\[
E\!\left[
\{X^{\prime}BX-\operatorname{tr}(B\Sigma_g)\}
\{X^{\prime}CX-\operatorname{tr}(C\Sigma_g)\}
\right]
=2\operatorname{tr}(B\Sigma_gC\Sigma_g),
\qquad X\sim N(0,\Sigma_g),
\]
therefore give, group by group,
\[
\operatorname{Var}_{\theta_0}(\Delta_n)\longrightarrow I(\theta_0),
\]
with exactly the two mean terms and two covariance terms in the definition of \(I(\theta_0)\).  Its assumed positive definiteness makes the limiting score nondegenerate.

Matrix inversion and the log determinant are smooth on the common spectral neighborhood.  Taylor expansion of each normal log density at chart displacement \(u/\sqrt n\), followed by summation over the finitely many groups, yields
\[
\log\frac{dP_{\theta_n(u)}}{dP_{\theta_0}}
=u^{\prime}\Delta_n-\frac12u^{\prime}I(\theta_0)u+r_n(u),
\qquad
\sup_{u\in\mathcal U}|r_n(u)|
\longrightarrow_{P_{\theta_0}}0.
\]
To justify the uniform remainder, the third derivatives of the normal log density are polynomials of fixed degree in a standardized Gaussian vector, with coefficients uniformly bounded on the spectral neighborhood.  Their sample averages are tight, while
\(\sup_{u\in\mathcal U}\lVert u/\sqrt n\rVert_2^3=O(n^{-3/2})\); summing \(O(n)\) observations gives \(o_{P_{\theta_0}}(1)\).  The same polynomial-Gaussian bound supplies a common exponential moment in a neighborhood of the origin.

For completeness, this exponential moment also upgrades the score central limit to total variation.  After standardization, each score summand is a polynomial of degree at most two in a fixed-dimensional standard Gaussian vector.  Its characteristic function is analytic near the origin.  Taylor expansion there, multiplication over the \(O(n)\) independent summands, and the preceding covariance limit give the characteristic function of \(N(0,I(\theta_0))\) uniformly on every fixed frequency ball.  Outside a fixed ball, positive group proportions and nondegeneracy of \(I(\theta_0)\) provide an integrable geometric majorant after a fixed number of convolutions; the explicit Gaussian-quadratic characteristic functions provide that majorant.  Fourier inversion and then enlargement of the ball imply
\[
\left\|
\mathcal L_{\theta_0}(\Delta_n)-N(0,I(\theta_0))
\right\|_{\mathrm{TV}}\longrightarrow0.
\]
The common exponential moment makes both the Taylor remainders and their likelihood-ratio products uniformly integrable.  Consequently the likelihood-ratio process converges in \(L^1(P_{\theta_0})\), uniformly on \(\mathcal U\), to
\[
\Lambda_u(\Delta)
=\exp\!\left\{
u^{\prime}\Delta-\frac12u^{\prime}I(\theta_0)u
\right\},
\qquad
\Delta\sim N(0,I(\theta_0)).
\]

We now construct the kernels.  Define an approximate law on the exact panel sample space by
\[
\frac{d\widetilde P_{n,u}}{dP_{\theta_0}}
=c_n(u)^{-1}
\exp\!\left\{
u^{\prime}\Delta_n-\frac12u^{\prime}I(\theta_0)u
\right\},
\]
where
\[
c_n(u)=E_{\theta_0}
\exp\!\left\{
u^{\prime}\Delta_n-\frac12u^{\prime}I(\theta_0)u
\right\}.
\]
Uniform \(L^1\) convergence gives
\[
\sup_{u\in\mathcal U}|c_n(u)-1|\longrightarrow0
\]
and
\[
\sup_{u\in\mathcal U}
\lVert\widetilde P_{n,u}-P_{\theta_n(u)}\rVert_{\mathrm{TV}}
\longrightarrow0.
\]
In the approximate experiment, \(\Delta_n\) is sufficient because every likelihood ratio is a measurable function of \(\Delta_n\).  The forward kernel reports \(\Delta_n\).  The reverse kernel draws from the center conditional law of the exact Gaussian sufficient statistic given \(\Delta_n\), and then regenerates the parameter-free ancillary residual rotations.  Such regular conditional laws exist because the sample and statistic spaces are Euclidean Borel spaces.

The center law of \(\Delta_n\) converges in total variation to \(N(0,I)\).  Multiplication by the uniformly integrable exponential tilts preserves this convergence uniformly for \(u\in\mathcal U\), so its law under \(\widetilde P_{n,u}\) converges uniformly to \(N(Iu,I)\).  A maximal common-part coupling between the center score law and \(N(0,I)\), followed by the same exponential tilt, supplies parameter-independent Markov kernels in both directions whose errors converge uniformly to zero.  Applying the linear map \(Z=I^{-1}\Delta\) yields \(Z\sim N(u,I^{-1})\).  The triangle inequality with the preceding approximation of \(P_{\theta_n(u)}\) proves that both deficiencies vanish uniformly on \(\mathcal U\).  Normal group likelihoods depend only on block means and residual scatter matrices, so these kernels can indeed be implemented from the sufficient statistics named in \(\mathrm{def:panel\mbox{-}transfer\mbox{-}handle}\); conditional simulation regenerates the discarded center-ancillary rotations.

It remains to transfer the band decision, including its target and unbounded-width issues.  By \(\mathrm{lem:local\mbox{-}risk\mbox{-}map}\),
\[
\beta_n:=
\sup_{u\in\mathcal U,\ t\in\mathcal D}
\left|
\sqrt n\{R_{P_{\theta_n(u)}}(v_n(t))
-R_{P_{\theta_0}}(v_0)\}
-\Psi_{\theta_0}(u,t)
\right|
\longrightarrow0.
\]
Define
\[
\Psi_n(u,t)=
\sqrt n\{R_{P_{\theta_n(u)}}(v_n(t))
-R_{P_{\theta_0}}(v_0)\}
\]
and
\[
I_n(t)=
\left[
\inf_{u\in\mathcal U}\Psi_n(u,t),\,
\sup_{u\in\mathcal U}\Psi_n(u,t)
\right].
\]
Compactness and uniform convergence to the continuous function \(\Psi_{\theta_0}\) imply a deterministic bound \(D<\infty\) on
\(\sup_t\operatorname{diam}I_n(t)\) for all sufficiently large \(n\).  Intersect every candidate section with \(I_n(t)\), returning a fixed point of \(I_n(t)\) if the intersection is empty.  Whenever the candidate covers a curve \(\Psi_n(u,\cdot)\), the intersection is nonempty and still covers it; width does not increase.  Thus the loss is bounded by \(D\).

If a transfer kernel has total-variation error \(\kappa_n\), it changes any simultaneous coverage probability by at most \(\kappa_n\) and any expected clipped loss by at most \(D\kappa_n\).  Padding both endpoints by \(\beta_n\) transfers between \(\Psi_n\) and \(\Psi_{\theta_0}\), adding at most \(2\beta_n\) to maximum width.  Finally, if the transferred rule has coverage \(q_{\mathrm{cov}}-\gamma_n\), mix it with the full deterministic envelope with probability
\[
\frac{\gamma_n}{\alpha+\gamma_n}.
\]
The resulting coverage is at least
\[
\left(1-\frac{\gamma_n}{\alpha+\gamma_n}\right)
(q_{\mathrm{cov}}-\gamma_n)
+\frac{\gamma_n}{\alpha+\gamma_n}
=q_{\mathrm{cov}},
\]
and the extra expected width is at most
\[
D\frac{\gamma_n}{\alpha+\gamma_n}.
\]
All three errors vanish.  Applying the construction in both kernel directions proves preservation of simultaneous coverage and expected maximum width and hence equality of the compact-local Gaussian-shift and observed-panel minimax bounded curve-diameter values.

%%% FIELD proof-local-panel-minimax
Let \(q_{\mathrm{cov}}=1-\alpha\).  By
\(\mathrm{thm:certified\mbox{-}gaussian\mbox{-}value}\), for each
\(\varepsilon_{\mathrm{opt}}>0\) there is a bounded finite padded-envelope Gaussian rule with simultaneous coverage \(q_{\mathrm{cov}}\) and worst expected maximum width at most
\[
L_\alpha(\theta_0,v_0;M,L_0)+\frac{\varepsilon_{\mathrm{opt}}}{3}.
\]
Apply the panel-to-Gaussian kernel from
\(\mathrm{thm:panel\mbox{-}experiment\mbox{-}transfer}\), pad its interval endpoints by the uniform target-approximation error, and mix with the full local envelope to repair the vanishing coverage error.  Centering at
\(R_{P_{\theta_0}}(v_0)\) and dividing by \(\sqrt n\) gives observed-panel intervals satisfying
\[
\liminf_{n\to\infty}\inf_{u\in\mathcal U}
P_{\theta_n(u)}
\left\{
R_{P_{\theta_n(u)}}(v_n(t))
\in\widetilde C_n^{(\varepsilon_{\mathrm{opt}})}(v_n(t))
\text{ for all }t\in\mathcal D
\right\}
\ge q_{\mathrm{cov}}
\]
and
\[
\limsup_{n\to\infty}\sup_{u\in\mathcal U}
E_{\theta_n(u)}
\left[
\sup_{t\in\mathcal D}
\sqrt n\,
\operatorname{diam}
\widetilde C_n^{(\varepsilon_{\mathrm{opt}})}(v_n(t))
\right]
\le
L_\alpha(\theta_0,v_0;M,L_0)+\varepsilon_{\mathrm{opt}}.
\]
The first display is exactly membership in
\(\mathfrak H^{\mathrm{panel}}_{1-\alpha}(\theta_0,v_0;M,L_0)\).

For the converse, suppose that a sequence in this honesty class had, along a subsequence, worst scaled expected maximum width at most
\(L_\alpha(\theta_0,v_0;M,L_0)-\epsilon\) for some \(\epsilon>0\).  Center and scale its sections and clip them to the deterministic local target envelope used in
\(\mathrm{thm:panel\mbox{-}experiment\mbox{-}transfer}\).  Coverage is preserved and loss is bounded.  The reverse kernel followed by target padding then produces Gaussian bands with coverage
\(q_{\mathrm{cov}}-o(1)\) and worst risk at most
\[
L_\alpha(\theta_0,v_0;M,L_0)-\epsilon+o(1).
\]
Mixing with the full Gaussian envelope with probability
\(o(1)/\{\alpha+o(1)\}\) restores coverage \(q_{\mathrm{cov}}\) at cost \(o(1)\).  For all sufficiently large indices this gives a member of
\(\mathfrak C_{1-\alpha}\) with worst risk strictly below
\(L_\alpha(\theta_0,v_0;M,L_0)\), contradicting the definition of that value.  This proves the lower bound.

Choose \(\varepsilon_k\downarrow0\) and increasing indices \(n_k\) so that, for every \(n\ge n_k\), the coverage deficit, transfer error, and excess risk for the \(k\)-th finite Gaussian rule are each at most \(\varepsilon_k\).  Use the \(k\)-th transferred and repaired rule for
\(n_k\le n<n_{k+1}\).  The resulting diagonal sequence is locally honest and has limsup loss at most \(L_\alpha\); the converse just proved makes its liminf loss at least \(L_\alpha\).  Hence its limit equals \(L_\alpha\), and its finite primal rules and dual bounds are exactly the certified programs in
\(\mathrm{thm:certified\mbox{-}gaussian\mbox{-}value}\).  No step uses the global projection band \(C_n\), so this argument asserts neither its optimality nor Gaussian-local attainment uniformly over the full mixing class.

%%% FIELD statement-global-honesty-local-converse
Let \(\{\widetilde C_n\}\) be randomized measurable observed-panel interval bands satisfying
\[
\liminf_{n\to\infty}\inf_{P\in\mathcal M_n}
P\!\left\{
R_P(v)\in\widetilde C_n(v)
\text{ for every }v\in\Delta_K
\right\}
\ge1-\alpha.
\]
For every compact Gaussian local path
\(\{\theta_n(u):u\in\mathcal U\}\subseteq\Theta_n^{\mathrm G}\), this sequence belongs, after restriction to \(v_n(\mathcal D)\), to
\(\mathfrak H^{\mathrm{panel}}_{1-\alpha}(\theta_0,v_0;M,L_0)\), and therefore
\[
\liminf_{n\to\infty}\sup_{u\in\mathcal U}
E_{P_{\theta_n(u)}}\!\left[
\sup_{t\in\mathcal D}\sqrt n\,
\operatorname{diam}\widetilde C_n(v_n(t))
\right]
\ge L_\alpha(\theta_0,v_0;M,L_0).
\]
This implication is one-sided: attainment within the Gaussian-local honesty class does not establish attainment by a sequence uniformly honest over \(\mathcal M_n\).

%%% FIELD proof-global-honesty-local-converse
Fix \(u\in\mathcal U\).  By
\(\mathrm{def:local\mbox{-}gaussian\mbox{-}parameter\mbox{-}space}\),
\(\theta_n(u)\in\Theta_n^{\mathrm G}\subseteq\Theta_n\).  By
\(\mathrm{lem:global\mbox{-}parameterization}\), the induced observed law
\(P_{\theta_n(u)}\) belongs to \(\mathcal M_n\).  Also
\(v_n(t)\in\Delta_K\) for all sufficiently large \(n\), uniformly over the compact tangent set \(\mathcal D\), by the defining tangent sign restrictions at zero coordinates of \(v_0\).  Hence the global simultaneous-coverage event
\[
\left\{
R_{P_{\theta_n(u)}}(v)
\in\widetilde C_n(v)
\text{ for every }v\in\Delta_K
\right\}
\]
is contained in the local event
\[
\left\{
R_{P_{\theta_n(u)}}(v_n(t))
\in\widetilde C_n(v_n(t))
\text{ for every }t\in\mathcal D
\right\}.
\]
Taking probabilities, then the infimum over \(u\), and finally the liminf proves membership in
\(\mathfrak H^{\mathrm{panel}}_{1-\alpha}(\theta_0,v_0;M,L_0)\).  The lower-bound half of
\(\mathrm{thm:local\mbox{-}panel\mbox{-}minimax}\) gives the displayed expected-width inequality.

The set of sequences uniformly honest over \(\mathcal M_n\) is therefore a subset of the locally honest Gaussian-path sequences.  Minimizing over the larger local class can only weakly decrease the attainable value.  Thus the lower bound transfers from the Gaussian subclass to every broader honesty requirement containing those paths, whereas an attaining member of the larger local class need not belong to the smaller globally honest class.  This proves the stated one-sidedness.

%%% FIELD construction-model-class-replacement
\(\mathcal M_n=\{P_n:P_n\text{ obeys every member property listed in }\texttt{by-member-properties}\}\), with the triangular-array index \(n\), all dimensions, and all fixed constants governed by those displayed properties.  No equality with, or additional restriction inherited from, any bibliographically cited comparison class is part of this definition.

%%% FIELD construction-local-gaussian-replacement
\(\Theta_n^{\mathrm G}=\{\theta\in\Theta_n:\mathsf Q_{\varepsilon}\text{ is the law of independent }N(0,\Sigma)\text{ vectors over all pre dates and induces independent }N(0,\Sigma_{DD})\text{ donor vectors over all observed post dates}\}\subseteq\Theta_n\), with the rank and spectral constraints inherited from \(\Theta_n\).  The displayed subset relation is part of the construction.

%%% FIELD construction-local-panel-honesty-replacement
\(\mathfrak H^{\mathrm{panel}}_{1-\alpha}(\theta_0,v_0;M,L_0)=\{\{\widetilde C_n\}_{n\ge1}:\widetilde C_n(v)\text{ is a randomized measurable interval based only on the observed Gaussian panel, and }\liminf_{n\to\infty}\inf_{u\in\mathcal U}\Pr_{P_{\theta_n(u)}}[R_{P_{\theta_n(u)}}(v_n(t))\in\widetilde C_n(v_n(t))\ \forall t\in\mathcal D]\ge1-\alpha\}\).  This definition imposes honesty only on the displayed compact iid-Gaussian local paths; it does not impose uniform honesty over \(\mathcal M_n\).
